\chapter{جمع‌بندی و چشم‌انداز}
\label{chap5}

\section{جمع‌بندی نتایج نمونه}

مدل پارامتری در فصل نخست در قالب یک مسئلهٔ مستقیم و یک مسئلهٔ معکوس صورت‌بندی شد. وجود و یکتایی جواب و وابستگی پیوسته به دادهٔ اولیه، مبنای تحلیلی محاسبات را فراهم کردند. مطالعهٔ پالایش گام نیز رفتار مرتبهٔ اول روش اویلر را برای مدل معیار نشان داد. در مطالعهٔ موردی ساختگی، خطای پارامتر، خطای برازش و حساسیت نسبت به گام زمانی به‌صورت جداگانه گزارش شدند.

\begin{table}[htbp]
  \centering
  \begin{tabular}{p{3.1cm}p{4.5cm}p{4.1cm}}
    \toprule
    مؤلفهٔ ارزیابی & شاهد نمونه & دامنهٔ نتیجه \\
    \midrule
    خوش‌تعریفی مدل
    & قضیهٔ وجود و یکتایی و کران گرونوالم
    & مسئلهٔ مستقیم و حساسیت به مقدار اولیه \\
    همگرایی عددی
    & خطای گام و مرتبهٔ تجربی
    & دقت گسسته‌سازی زمانی \\
    برازش پارامتر
    & خطای پارامتر و $\operatorname{RMSE}$
    & سناریوی دادهٔ ساختگی \\
    بازتولیدپذیری
    & فرادادهٔ کامل اجرای محاسباتی
    & اجرای محاسباتی مشخص \\
    \bottomrule
  \end{tabular}
  \caption{پیوند میان پرسش، شاهد و دامنهٔ اعتبار نتیجه}
  \label{tab:evidence-map}
\end{table}

\section{محدودیت‌ها}

مدل معیار یک‌بعدی و داده‌ها ساختگی‌اند؛ ساختار خطا نیز به نویز مستقل با واریانس محدود شده است. ازاین‌رو، نتایج نمونه دربارهٔ سامانه‌های غیرخطی، داده‌های هم‌بسته، پارامترهای چندگانه یا عدم‌قطعیت ساختاری مدل ادعایی ندارند. همچنین یک اجرای منفرد برای استنباط آماری کافی نیست و جایگزین تحلیل تکرارپذیری و فاصلهٔ اطمینان نمی‌شود.

\section{مسیرهای توسعه}

گسترش مدل به دستگاه‌های غیرخطی یا تصادفی، تحلیل شناسایی‌پذیری چندپارامتری، استفاده از روش‌های تطبیقی و مرتبه‌بالاتر، و کمّی‌سازی عدم‌قطعیت از مسیرهای طبیعی ادامهٔ مطالعه‌اند. در هر توسعه، نتیجهٔ تحلیلی، آزمون همگرایی عددی و شاهد مبتنی بر داده دامنه‌های اعتبار جداگانه دارند و در کنار یکدیگر تصویر کامل‌تری از مدل فراهم می‌کنند.
